\documentclass[11pt,a4paper]{article}

% ── Packages ──────────────────────────────────────────────────────────────────
\usepackage[margin=2.4cm]{geometry}
\usepackage{amsmath,amssymb}
\usepackage{graphicx}
\usepackage{booktabs}
\usepackage{array}
\usepackage{xcolor}
\usepackage{listings}
\usepackage{enumitem}
\usepackage{caption}
\usepackage{subcaption}
\usepackage[hidelinks]{hyperref}
\usepackage{microtype}

% ── Code listing style ────────────────────────────────────────────────────────
\definecolor{codebg}{rgb}{0.97,0.97,0.97}
\definecolor{codekw}{rgb}{0.10,0.10,0.55}
\definecolor{codecom}{rgb}{0.30,0.50,0.30}
\definecolor{codestr}{rgb}{0.60,0.20,0.10}
\lstdefinestyle{py}{
  language=Python,
  backgroundcolor=\color{codebg},
  basicstyle=\ttfamily\footnotesize,
  keywordstyle=\color{codekw}\bfseries,
  commentstyle=\color{codecom}\itshape,
  stringstyle=\color{codestr},
  numbers=left, numberstyle=\tiny\color{gray}, numbersep=6pt,
  frame=single, rulecolor=\color{gray!40},
  breaklines=true, showstringspaces=false,
  tabsize=2, columns=fullflexible, keepspaces=true,
}
\lstset{style=py}

\newcommand{\file}[1]{\texttt{#1}}
\newcommand{\code}[1]{\texttt{#1}}
% Wrapping filename for narrow table cells (footnotesize typewriter).
\newcommand{\fn}[1]{{\footnotesize\ttfamily #1}}
% Ragged-right paragraph column.
\newcolumntype{L}[1]{>{\raggedright\arraybackslash}p{#1}}

% ── Title ─────────────────────────────────────────────────────────────────────
\title{\textbf{From Optimal Control to Diffusion Policy}\\[2pt]
\large A Trajectory-Optimization $\rightarrow$ TVLQR $\rightarrow$
Behavioural-Cloning Pipeline for Double-Pendulum Swing-Up}
\author{Group Project TN2 --- Technical Documentation}
\date{\today}

\begin{document}
\maketitle
\clearpage

\tableofcontents
\bigskip
\hrule

\clearpage
% ══════════════════════════════════════════════════════════════════════════════
\section{Introduction}
The goal of this project is to compare the performance of 2 different learning-
based control methods. The first method is a behavioural cloning approach, 
where we use a dataset of expert demonstrations to train a policy 
that mimics the expert's actions. The second method is a diffusion policy approach,
where we use a diffusion model to learn a policy that can generate actions based on
the current state of the system. We will evaluate the performance of both methods
in the \textbf{double pendulum swingup task} and the \textbf{stacking problem}.

\subsection{Behavioural Cloning}

\subsection{Diffusion}
In order to understand the diffusion policy approach, we first need to understand
the concept of diffusion models. Diffusion models are a class of generative models
that learn to generate structured data (in our case, action sequences) 
by iteratively denoising a sample of noise.

\subsection*{Parameter Setup}
For the diffusion model, we first need to define the time parameters of the 
noising and denoising process.
\begin{itemize}[leftmargin=1.5em]
  \item \textbf{Time Horizon} $T$ is the number of time steps in the diffusion process.
  \item \textbf{$\beta$ Schedule} $\beta_t$ is a set of hyperparameters that control
  how quickly the noise is added to the data during the forward diffusion process.
  It is a monotonic increasing sequence of values that determine the variance of the noise
  \item \textbf{Noise Schedule} $\alpha_t$ = $1 - \beta_t$ is a set of hyperparameters 
  that control how quickly the noise is removed.
  \item \textbf{Cumulative Noise Schedule} $\bar{\alpha}_t$ = $\prod_{s=1}^{t} \alpha_s$ 
  \item 
\end{itemize}


\end{document}
